\documentclass[10pt,a4paper]{amsart}
\usepackage[margin=19mm]{geometry}
\usepackage{amsmath,amssymb,mathtools}
\usepackage[scr=rsfs,cal=euler]{mathalfa}
\usepackage{xfrac}
\usepackage[unicode,hidelinks,bookmarks=false]{hyperref}
\newcommand{\ie}{i.e.\ }
\newcommand{\cf}{cf.\ }
\newcommand{\resp}{respectively\ }
\newcommand{\Subsection}{\subsection}
\providecommand{\nobreakdash}{\nobreak\hbox{-}}
\setlength{\emergencystretch}{2em}
\multlinegap0pt
\usepackage{bm}
\usepackage{bbm}
\usepackage{dsfont}
\usepackage{xspace}
\usepackage{cancel}
\usepackage{mathtools}
\usepackage{amsthm}
\makeatletter
\@ifundefined{lemma}{\newtheorem{lemma}{Lemma}}{}
\@ifundefined{theorem}{\newtheorem{theorem}{Theorem}}{}
\makeatother
\title[On the Hausdorff dimension of graph of random vector-valued ]{On the Hausdorff dimension of graph of random vector-valued Weierstrass function}
\author{Jun Jason Luo, Zi-Rui Zhang}
\date{Source: arXiv:2604.13913v2 (CC BY 4.0)}
\begin{document}
\maketitle

\noindent\emph{Source and licence.} On the Hausdorff dimension of graph of random vector-valued Weierstrass function. Version arXiv:2604.13913v2 (CC BY 4.0), distributed under the Creative Commons Attribution 4.0 International licence, as recorded in the arXiv licence metadata for this exact version and shown on the abstract page https://arxiv.org/abs/2604.13913v2. Full rights notice: licenses/arxiv-2604.13913-CC-BY-4.0.txt.\par\smallskip

\noindent\textit{Editorial scope.} Counted unit: the Theorem whose source label is \texttt{None} in the article (source0.tex lines 352--355, source label \texttt{None}), delivered together with the article's own complete original proof of it (source0.tex lines 358--478, one \texttt{proof} environment of 121 lines). No mathematics is written, repaired or re-derived: statement and proof are the article's own text line for line. Required context retained uncounted: eq-phi-n (lines 168--170); eq- vector-valued funct. (lines 171--173); lem-lowerbound (lines 207--209); thm2.1 (lines 233--235). Every definition, display and label of those lines is retained unchanged. Omitted from this package and not redistributed: 1--167, 174--206, 210--232, 236--351, 356--357, 479--583. Nothing inside a retained excerpt is omitted or altered: every retained body is delivered exactly as the article's lines are written, so no omission receipt of any kind accompanies this package. Citations: the retained excerpts carry no citation command at all, so no bibliography entry of the article is transcribed and no reference is redistributed. Reference adaptation: none was needed and none was made; every reference inside the delivered unit resolves inside it, so no unresolved reference or citation is produced. Printed slips kept literal: no printed word, symbol, hypothesis or number of the counted statement or of its proof was altered. Layout and class: the article's own class is \texttt{amsart} with packages including amsmath, amssymb, amsthm, mathrsfs, amsfonts, hyperref, enumitem, graphicx, upgreek, mathtools, tikz; the wrapper is the lane's pinned \texttt{amsart} editorial wrapper at 10pt on A4, which restates the theorem-like environments the retained text needs and adds the article's own macro definitions that the retained text needs (none), with extra wrapper packages (bm, bbm, dsfont, xspace, cancel, mathtools). The delivered numbering is the unit's own. Assets: no retained span contains a figure environment, a graphics-inclusion command, a PDF-page inclusion, a table or a TikZ picture; no image file is copied into this package, the package declares no asset dependency and no image file is redistributed with it. Licence: version arXiv:2604.13913v2 of this article is distributed under the Creative Commons Attribution 4.0 International licence, as recorded in the arXiv licence metadata for this exact version and shown on the abstract page https://arxiv.org/abs/2604.13913v2. A full rights notice is delivered at licenses/arxiv-2604.13913-CC-BY-4.0.txt. Characters: every retained excerpt is pure ASCII, so the delivered engine, XeLaTeX with its default Latin Modern text font, reports no missing character.
\begin{equation}\label{eq-phi-n}
	\varphi_n:=\varphi_n(x,y)=2b^{-\beta n}\sin(\pi b^n(x-y)). 
\end{equation} The sum-to-product formulas in trigonometry imply that

\begin{equation}\label{eq- vector-valued funct.}
	f_{\Theta,\Lambda}(x) -f_{\Theta,\Lambda}(y)=(a_1(x,y), a_2(x,y))
\end{equation}

	\begin{lemma}\label{lem-lowerbound}
		$\dim_H G(f_{\Theta,\Lambda})\le \min\left\{\frac{1}{\beta}, \, 3-2\beta\right\}$ a.s.
	\end{lemma}

\begin{theorem}\label{thm2.1}
 	$\dim_H f_{\Theta,\Lambda}([0,1]) = \min\{2, \frac{1}{\beta}\}$ a.s.
 \end{theorem}

 \begin{theorem} 
 	The Hausdorff dimension of the graph $f_{\Theta,\Lambda}$ satisfies
 	$$\dim_H G(f_{\Theta,\Lambda}) =\dim_B G(f_{\Theta,\Lambda}) = \min\left\{\frac{1}{\beta}, \, 3 -2\beta\right\} \quad \text{a.s.}$$
 \end{theorem}

\begin{proof} 
	Define the projection  $P: G(f_{\Theta,\Lambda}) \to f_{\Theta,\Lambda}([0,1])$ from the graph onto the image by 
	$$P(x, f_{\Theta,\Lambda}(x))= f_{\Theta,\Lambda}(x), \quad x\in [0,1].$$
	Trivially $P$ is a Lipschitz function.  Hence $\dim_H G(f_{\Theta,\Lambda})\ge \dim_H f_{\Theta,\Lambda}([0,1]).$
	
    If $\beta\in [\frac12, 1]$, 	it follows from  Lemma \ref{lem-lowerbound} and Theorem \ref{thm2.1} that
	$$ \frac {1}{\beta} \ge \dim_H G(f_{\Theta,\Lambda})\ge \dim_H f_{\Theta,\Lambda}([0,1])=\frac{1}{\beta} \quad \text{a.s.}$$
	That proves  $\dim_H G(f_{\Theta,\Lambda}) = \frac {1}{\beta}$ a.s.
	
	 If $\beta\in (0, \frac12)$,  we let   \( \mu:= \mu_{\Theta,\Lambda} \) be the induced measure on $\mathbb{R}^3$ through the Lebesgue measure  \( {\mathcal L} \) and $f_{\Theta,\Lambda}$ in the following way: for any Borel set \( V \subset \mathbb{R}^3 \), define
	\begin{equation*}
		\mu (V) = {\mathcal L}\bigl(\{t \in [0,1] : (t, f_{\Theta,\Lambda}(t)) \in V\}\bigr).
	\end{equation*} 
	The measure $\mu$ is supported on the graph  $G(f_{\Theta,\Lambda})$.  
	
	For $s>0$, the \(s\)-energy of \( \mu \)  is given by
	\begin{eqnarray*}
			I_s(\mu)
			&=& \int_{\mathbb{R}^3} \int_{\mathbb{R}^3}\frac{d\mu(w)d\mu(z)}{|w-z|^s} \\
			&=& \int_{[0,1]^2} \frac{d(x,y)}{\bigl((x-y)^2 + |f_{\Theta,\Lambda}(x)-f_{\Theta,\Lambda}(y)|^2\bigr)^{s/2}}.
	\end{eqnarray*}
	
	We will prove that for any \( s\in (2, 3-2\beta) \), the \(s\)-energy \( I_s(\mu) \) is finite for almost every pair of phase sequences \( (\Theta,\Lambda) \in H \times H \). This implies that the Hausdorff dimension of the graph of \( f_{\Theta,\Lambda} \) is at least \( s \). By letting \( s \to 3-2\beta \), we conclude that
	\[
	\dim_H G(f_{\Theta,\Lambda}) \ge 3-2\beta
	\]
	holds for almost every \( (\Theta,\Lambda) \). Hence the desired result follows by Lemma \ref{lem-lowerbound}.
	
	In order to show $I_s(\mu)<\infty $ a.s., it suffices to prove that the expectation \( \mathbb{E}(I_s(\mu)) < \infty \). By Fubini-Tonelli theorem, we have
	\begin{equation*}
		\mathbb{E}(I_s(\mu))= \int_{[0,1]^2} \mathbb{E}\left(\frac{1}{\bigl((x-y)^2 + |f_{\Theta,\Lambda}(x)-f_{\Theta,\Lambda}(y)|^2\bigr)^{s/2}}\right) \, d(x,y).
	\end{equation*}
	
	We claim that there exists $c_0>0$ such that for \( 0 < |x-y| \le \frac{1}{2b^2} \),  
	$$\mathbb{E}\left(\frac{1}{\bigl((x-y)^2 + |f_{\Theta,\Lambda}(x)-f_{\Theta,\Lambda}(y)|^2\bigr)^{s/2}}\right)\le c_0|x-y|^{2-s-2\beta}.$$
	Since \( 2 < s < 3-2\beta \), it concludes that $\mathbb{E}(I_s(\mu))<\infty$.
	
	Fix \( x,y \in [0,1] \) satisfying \( 0 < |x-y| \le \frac{1}{2b^2}\), and define
	$$X = |f_{\Theta,\Lambda}(x)-f_{\Theta,\Lambda}(y)|^2.$$
	Regarding \( X \) as a random variable depending on \( (\Theta,\Lambda) \), it can be verified that \( X \) has a bounded probability density function \( h \). Then
	\begin{eqnarray*}
			&& \mathbb{E}\left(\frac{1}{\bigl((x-y)^2 + |f_{\Theta,\Lambda}(x)-f_{\Theta,\Lambda}(y)|^2\bigr)^{s/2}}\right) \\
			&= &\int_0^\infty \frac{h(t)\,dt}{\bigl((x-y)^2 + t\bigr)^{s/2}} \\
			&= &\int_0^\infty \frac{h\bigl(|x-y|^2 \tau\bigr) |x-y|^2 d\tau}{|x-y|^s (1+\tau)^{s/2}} \\
			&\le &\Bigl(\sup_{t\ge 0} h(t)\Bigr) \, |x-y|^{2-s} \int_0^\infty \frac{d\tau}{(1+\tau)^{s/2}}.
	\end{eqnarray*}
	 Hence to prove the claim, we need to show that there exists a constant \( c_1>0 \), independent of \( x,y \), such that
	\begin{equation}\label{eq-claim}
	h(t) \le c_1|x-y|^{-2\beta}.
	\end{equation}
	
	By using the notation in \eqref{eq-phi-n} and \eqref{eq- vector-valued funct.}, $X$ can be written as
	$$X = (a_1(x,y))^2 +(a_2(x,y))^2 := X_1^2 + X_2^2,$$
	where    
	   $$X_1 = \sum_{n=0}^\infty  \varphi_n \sin\left(\pi\left(b^n x + b^n y +2\theta_n\right)\right)$$ 
	and 
	   $$X_2 = \sum_{n=0}^\infty  \varphi_n \cos\left(\pi\left(b^n x + b^n y +2\lambda_n\right)\right).$$
   Since \( \theta_n \sim \mathrm{Uniform}[0,1] \) are i.i.d., each \( \varphi_n \sin\left(\pi\left(b^n x + b^n y +2\theta_n\right)\right) \) has density function
	\begin{equation*}
		h_n(z) =
		\begin{cases}
			\dfrac{1}{\pi\sqrt{{\varphi_n}^2-z^2}} & |z| < |\varphi_n|, \\
			0 & |z| \ge |\varphi_n|.
		\end{cases}
	\end{equation*}
 The density function of $X_1$ is the infinite convolution 
 $$h_{X_1} = h_0 * h_1 * h_2 * \cdots.$$ 
 Since the supremum of a probability density is nonincreasing under convolution, it follows that	
	\[
	\sup_t h_{X_1}(t) \le \sup_t \bigl(h_{k-2} * h_{k-1} * h_k\bigr)(t)
	\]
	for any fixed \( k \ge 2 \). 
	
	Choose an integer \( k \) such that $\frac{1}{2b^{k+1}} < |x-y| < \frac{1}{2b^k}$. 	Then
	\begin{equation*}
		\frac{\pi}{2b^3} < \Bigl|\pi b^{k-2}(y-x)\Bigr|\le \Bigl|\pi b^k(y-x)\Bigr| \le \frac{\pi}{2}.
	\end{equation*}
	So for \( n = k-2,k-1,k \),
	\begin{equation*}
		|\varphi_n|=|2b^{-\beta n}\sin(\pi b^n(x-y))| > 2\sin\Bigl(\frac{\pi}{2b^3}\Bigr) b^{-k\beta}
		> 2^{1+\beta}\sin\Bigl(\frac{\pi}{2b^3}\Bigr) |x-y|^\beta.
	\end{equation*}
	
	Let \( \|\cdot\|_p \) denote the \( L^p \)-norm. For \( p=3/2 \),
	\begin{eqnarray*}
		\|h_n\|_{\frac32}
			&=& \left( \int_{-|\varphi_n|}^{|\varphi_n|} \frac{dz}{\bigl(\pi\sqrt{{\varphi_n}^2-z^2}\bigr)^{\frac32}} \right)^{\frac23} \\
			&=& \pi^{-1} \left( \int_{-|\varphi_n|}^{|\varphi_n|} ({\varphi_n}^2-z^2)^{-\frac34}dz \right)^{\frac23} \\
			&=& \pi^{-1} I^{\frac23} |\varphi_n|^{-\frac13},
	\end{eqnarray*}
	where \( I = \int_{-1}^1 (1-u^2)^{-\frac34}du < \infty \) is a universal constant. Thus for \( n=k-2,k-1,k \),
	\[
	\|h_n\|_{\frac32} \le c_2 |x-y|^{-\frac{\beta}{3}}
	\]
	where \( c_2 \) depends only on \( b \). 	By Young's inequality and H\"older's inequality,
	\begin{eqnarray*}
		\bigl(h_{k-2} * h_{k-1} * h_k\bigr)(t)
			&\le& \|h_{k-2}\|_{\frac32} \|h_{k-1} * h_k\|_3 \\
			&\le&  \|h_{k-2}\|_{\frac32} \|h_{k-1}\|_{\frac32}\|h_k\|_{\frac32} \\
			& \le& c_3 |x-y|^{-\beta}
	\end{eqnarray*} 
	where $c_3 :=(c_2)^3$. Hence $h_{X_i}(t) \le c_3 |x-y|^{-\beta}$. The same bound holds for  $h_{X_2}$, the density function of $X_2$.
	
	Finally, we compute the density function of $X= X_1^2 + X_2^2$. For \( i=1,2 \), the density of \( X_i^2 \) is
	\begin{equation*}
		h_{X_i^2}(t) =
		\begin{cases}
			\dfrac{h_{X_i}(\sqrt{t})+h_{X_i}(-\sqrt{t})}{2\sqrt{t}}, & t\ge 0, \\
			0, & t<0.
		\end{cases}
	\end{equation*}
	Since \( X_1 \) and \( X_2 \) are independent, the density function of \( X \) satisfies
	\begin{eqnarray*}
			h(t) &=& \bigl(h_{X_1^2} * h_{X_2^2}\bigr)(t) \\
			&=& \int_0^t \frac{h_{X_1}(\sqrt{t-s})+h_{X_1}(-\sqrt{t-s})}{2\sqrt{t-s}}
			\cdot \frac{h_{X_2}(\sqrt{s})+h_{X_2}(-\sqrt{s})}{2\sqrt{s}} ds \\
			&\le& ({c_3})^2|x-y|^{-2\beta} \int_0^t \frac{ds}{\sqrt{s(t-s)}} \\
			&=& ({c_3})^2\pi |x-y|^{-2\beta}.
	\end{eqnarray*}
By letting $c_1 :=\pi ({c_3})^2$, we obtain the inequality \eqref{eq-claim}, hence prove the theorem.
\end{proof}

\end{document}
